% GENERATED FILE. Edit canonical lessons, then run npm run generate:books.
% canonical-source-hash: fnv1a64:fbe00dffa2bc70cb
% canonical-lessons: ES-C714-carino, ES-C714-carinoso, ES-C714-ternura, ES-C714-tierno, ES-C714-adorar

\chapter{Cariño y ternura}
\label{ch:es-carino-y-ternura}
\hlchaptermodality{714}

\begin{chapteropening}
\textbf{By the end of this chapter:} \emph{I can show affection and talk about tenderness in a family.}

Complete the last lesson of chapter 714: I can show affection and talk about tenderness in a family.
\end{chapteropening}

\section[el cariño]{el cariño — affection; darling}

\label{lesson:ES-C714-carino}



\begin{quote}
Before the new one: say the Spanish for to get bored, then the Spanish for boredom.
\end{quote}

\subsection*{You'll want to know: el cariño}
\textbf{el cariño} — \textquotedblleft{}affection; darling\textquotedblright{}.

\begin{quote}
Le tengo mucho cariño a mi abuela.
\end{quote}

\begin{quote}
\emph{I am very fond of my grandmother.}
\end{quote}

\emph{Tener cariño a alguien} — to be fond of someone. \emph{Cariño} is also a pet name: darling.

\begin{cousinweb}
\emph{cariño} was borrowed from Galician-Portuguese \emph{cariño} (EtymDB). Further back the origin is not settled; it is often linked to Latin \emph{carēre} \textquotedblleft{}to lack\textquotedblright{}, as longing for someone absent.
\end{cousinweb}

\begin{grammarlens}[title={What you've built}]
One more word for telling, explaining and describing at B1.
\end{grammarlens}

\subsection*{Your turn}
\begin{itemize}
\raggedright
  \item \emph{Say it:} \emph{el cariño}
  \item \emph{Say it:} \emph{el cariño}, once more
  \item \emph{Say it:} a sentence of your own with \emph{el cariño}
  \item \emph{Recall:} say the Spanish for to get bored, then the Spanish for boredom, then say \emph{el cariño} again
\end{itemize}

\begin{tcolorbox}[breakable,colback=teal!4,colframe=teal!35!black,title={Before you move on}]
What does el cariño mean? (\textquotedblleft{}Affection; darling\textquotedblright{}.) Say it once more.
\end{tcolorbox}
\section[cariñoso]{cariñoso — affectionate, loving}

\label{lesson:ES-C714-carinoso}



\begin{quote}
Before the new one: say the Spanish for boredom, then the Spanish for affection.
\end{quote}

\subsection*{You'll want to know: cariñoso}
\textbf{cariñoso} — \textquotedblleft{}affectionate, loving\textquotedblright{}.

\begin{quote}
Mi perro es muy cariñoso.
\end{quote}

\begin{quote}
\emph{My dog is very affectionate.}
\end{quote}

Used with \emph{ser} for character, with \emph{estar} for how someone is acting now.

\begin{cousinweb}
\emph{cariñoso} is built on \emph{cariño} \textquotedblleft{}affection\textquotedblright{} with the ending \emph{-oso}; \emph{cariño} was borrowed from Galician-Portuguese (EtymDB).
\end{cousinweb}

\begin{grammarlens}[title={What you've built}]
One more word for telling, explaining and describing at B1.
\end{grammarlens}

\subsection*{Your turn}
\begin{itemize}
\raggedright
  \item \emph{Say it:} \emph{cariñoso}
  \item \emph{Say it:} \emph{cariñoso}, once more
  \item \emph{Say it:} a sentence of your own with \emph{cariñoso}
  \item \emph{Recall:} say the Spanish for boredom, then the Spanish for affection, then say \emph{cariñoso} again
\end{itemize}

\begin{tcolorbox}[breakable,colback=teal!4,colframe=teal!35!black,title={Before you move on}]
What does cariñoso mean? (\textquotedblleft{}Affectionate, loving\textquotedblright{}.) Say it once more.
\end{tcolorbox}
\section[la ternura]{la ternura — tenderness}

\label{lesson:ES-C714-ternura}



\begin{quote}
Before the new one: say the Spanish for affection, then the Spanish for affectionate.
\end{quote}

\subsection*{You'll want to know: la ternura}
\textbf{la ternura} — \textquotedblleft{}tenderness\textquotedblright{}.

\begin{quote}
Mira a su hijo con ternura.
\end{quote}

\begin{quote}
\emph{She looks at her son with tenderness.}
\end{quote}

\emph{Me da ternura} — it melts my heart.

\begin{cousinweb}
\emph{ternura} is built on \emph{tierno} \textquotedblleft{}tender\textquotedblright{} with the ending \emph{-ura}; \emph{tierno} is inherited from Latin \emph{tener} \textquotedblleft{}soft, delicate\textquotedblright{} (EtymDB).
\end{cousinweb}

\begin{grammarlens}[title={What you've built}]
One more word for telling, explaining and describing at B1.
\end{grammarlens}

\subsection*{Your turn}
\begin{itemize}
\raggedright
  \item \emph{Say it:} \emph{la ternura}
  \item \emph{Say it:} \emph{la ternura}, once more
  \item \emph{Say it:} a sentence of your own with \emph{la ternura}
  \item \emph{Recall:} say the Spanish for affection, then the Spanish for affectionate, then say \emph{la ternura} again
\end{itemize}

\begin{tcolorbox}[breakable,colback=teal!4,colframe=teal!35!black,title={Before you move on}]
What does la ternura mean? (\textquotedblleft{}Tenderness\textquotedblright{}.) Say it once more.
\end{tcolorbox}
\section[tierno]{tierno — tender, soft; sweet}

\label{lesson:ES-C714-tierno}



\begin{quote}
Before the new one: say the Spanish for affectionate, then the Spanish for tenderness.
\end{quote}

\subsection*{You'll want to know: tierno}
\textbf{tierno} — \textquotedblleft{}tender, soft; sweet\textquotedblright{}.

\begin{quote}
La carne está muy tierna.
\end{quote}

\begin{quote}
\emph{The meat is very tender.}
\end{quote}

Used for food (tender) and for people (sweet, gentle).

\begin{cousinweb}
\emph{tierno} is inherited from Latin \emph{tener}, \emph{tenerum} \textquotedblleft{}soft, delicate\textquotedblright{} (EtymDB); the stressed \emph{e} became \emph{ie}, as in \emph{tiempo}.
\end{cousinweb}

\begin{grammarlens}[title={What you've built}]
One more word for telling, explaining and describing at B1.
\end{grammarlens}

\subsection*{Your turn}
\begin{itemize}
\raggedright
  \item \emph{Say it:} \emph{tierno}
  \item \emph{Say it:} \emph{tierno}, once more
  \item \emph{Say it:} a sentence of your own with \emph{tierno}
  \item \emph{Recall:} say the Spanish for affectionate, then the Spanish for tenderness, then say \emph{tierno} again
\end{itemize}

\begin{tcolorbox}[breakable,colback=teal!4,colframe=teal!35!black,title={Before you move on}]
What does tierno mean? (\textquotedblleft{}Tender, soft; sweet\textquotedblright{}.) Say it once more.
\end{tcolorbox}
\section[adorar]{adorar — to adore, to worship}

\label{lesson:ES-C714-adorar}



\begin{quote}
Before the new one: say the Spanish for tenderness, then the Spanish for tender.
\end{quote}

\subsection*{You'll want to know: adorar}
\textbf{adorar} — \textquotedblleft{}to adore, to worship\textquotedblright{}.

\begin{quote}
Adoro el chocolate.
\end{quote}

\begin{quote}
\emph{I adore chocolate.}
\end{quote}

Stronger than \emph{encantar}; it also means to worship.

\begin{cousinweb}
\emph{adorar} was borrowed from Latin \emph{adōrāre} \textquotedblleft{}to address, to pray to\textquotedblright{} (EtymDB), built on \emph{ad} \textquotedblleft{}to\textquotedblright{} and \emph{ōrāre} \textquotedblleft{}to speak, to pray\textquotedblright{}.
\end{cousinweb}

\begin{grammarlens}[title={What you've built}]
One more word for telling, explaining and describing at B1.
\end{grammarlens}

\subsection*{Your turn}
\begin{itemize}
\raggedright
  \item \emph{Say it:} \emph{adorar}
  \item \emph{Say it:} \emph{adorar}, once more
  \item \emph{Say it:} a sentence of your own with \emph{adorar}
  \item \emph{Recall:} say the Spanish for tenderness, then the Spanish for tender, then say \emph{adorar} again
\end{itemize}

\begin{tcolorbox}[breakable,colback=teal!4,colframe=teal!35!black,title={Before you move on}]
What does adorar mean? (\textquotedblleft{}To adore, to worship\textquotedblright{}.) Say it once more.
\end{tcolorbox}
